%% Chương 3: Phương pháp đề xuất — Hybrid
\clearpage
\phantomsection
\section*{CHƯƠNG 3. PHƯƠNG PHÁP ĐỀ XUẤT: THUẬT TOÁN HYBRID}
\addcontentsline{toc}{section}{\numberline{}CHƯƠNG 3. PHƯƠNG PHÁP ĐỀ XUẤT: THUẬT TOÁN HYBRID}
\setcounter{section}{3}
\setcounter{subsection}{0}
\setcounter{figure}{0}
\setcounter{table}{0}
\setcounter{equation}{0}

\subsection{Động lực thiết kế}

\subsubsection{Phân tích yêu cầu}

Một thuật toán mã hoá chọn lọc DC H.264 hiệu quả cần đáp ứng đồng thời ba yêu cầu
thường xung đột nhau:

\begin{enumerate}
  \item \textbf{Bảo mật mạnh:} NPCR $> 99.6\%$, Entropy $> 7.9$, SSIM DC $\approx 0$.
        Điều này đòi hỏi cả \emph{keystream variation per NALU} và \emph{intra-NALU diffusion}.

  \item \textbf{Hiệu năng chấp nhận được:} DC encrypt time phải đủ thấp để
        không trở thành bottleneck trong pipeline (đặc biệt với S1: 19.037 NALU).
        Mục tiêu: $< 100$ ms tổng cho S1.

  \item \textbf{Byte-exact decryptability:} Giải mã phải tái tạo DC bytes gốc
        chính xác từng bit. BER = 0 là invariant cứng không thể vi phạm.
\end{enumerate}

\subsubsection{Tại sao cần per-NALU keystream variation}

Đây là yêu cầu thiết kế quan trọng nhất. Xét kịch bản: kẻ tấn công có thể
quan sát nhiều NALU encrypted trong cùng file. Nếu keystream giống nhau:

\begin{equation}
C_1 \oplus C_2 = (P_1 \oplus ks) \oplus (P_2 \oplus ks) = P_1 \oplus P_2
\end{equation}

Kẻ tấn công biết $C_1 \oplus C_2 = P_1 \oplus P_2$ mà không cần biết $ks$.
Nếu biết một trong hai plaintext (thông qua phân tích thống kê video), hắn suy ra
plaintext còn lại. Đây là tấn công ``many-time pad'' --- vi phạm IND-CPA.

Giải pháp: thay đổi keystream cho mỗi NALU bằng cách sử dụng \texttt{nalu\_index}
như một ``tweak'' --- giá trị công khai (không bí mật) nhưng khác nhau mỗi NALU,
đảm bảo keystream khác nhau ngay cả với cùng khóa K.

\subsubsection{Tại sao cần chained diffusion}

Chỉ có per-NALU keystream variation chưa đủ. Xét tấn công flip 1 bit:
kẻ tấn công flip byte $DC'[0] = DC[0] \oplus 1$ và gửi cho server.
Nếu chỉ dùng XOR đơn (không diffusion):

\begin{equation}
C'[i] = \text{permuted}'[i] \oplus ks[i]
\end{equation}

Chỉ $C'[0]$ thay đổi (do chỉ $\text{permuted}'[0]$ thay đổi). Vậy NPCR $= 1/25 = 4\%$.

Cơ chế chained diffusion giải quyết điều này: khi $C[0]$ thay đổi, nó làm $C[1]$
thay đổi (vì $C[1]$ phụ thuộc $C[0]$), từ đó $C[2], \ldots, C[24]$ đều thay đổi.
Kết quả: 1 bit thay đổi plaintext → 25 bytes thay đổi ciphertext → NPCR $\approx 99.6\%$.

\subsubsection{Tại sao cần Arnold permutation}

Arnold permutation phục vụ hai mục đích:

\begin{enumerate}
  \item \textbf{Phá vỡ cấu trúc cục bộ:} DC bytes liền kề trong NALU thường
        tương quan (cùng macroblock hoặc macroblock lân cận). Hoán vị Arnold
        trộn các bytes này trước khi XOR, làm mất tương quan cục bộ.

  \item \textbf{Tăng thành phần keyspace:} Tham số $(P, Q)$ của Arnold map
        là thành phần bổ sung của khóa, mặc dù keyspace nhỏ ($\approx 2^5$).

  \item \textbf{Tăng UACI:} Hoán vị làm thay đổi vị trí bytes trước diffusion,
        lý thuyết tăng UACI bằng cách phân phối đều hiệu ứng thay đổi.
        Tuy nhiên, Arnold fixed point ở một số NALUs làm UACI chưa đạt $33.46\%$
        (phân tích chi tiết ở Chương 6).
\end{enumerate}

\subsection{Kiến trúc tổng thể}

Thuật toán hybrid kết hợp ba tầng bảo mật theo thứ tự:

\begin{equation}
\text{DC bytes} \xrightarrow{\text{PLCM keystream}} ks[0..24]
\xrightarrow{\text{Arnold}} \text{permuted}[0..24]
\xrightarrow{\text{XOR + diffusion}} C[0..24]
\end{equation}

\begin{figure}[ht]
\centering
%% TikZ: Hybrid algorithm flowchart
\begin{tikzpicture}[
  node distance=1.1cm and 1.5cm,
  block/.style={draw, rectangle, rounded corners=3pt,
                minimum width=3.2cm, minimum height=0.7cm,
                text centered, font=\small},
  io/.style={block, fill=gray!15},
  proc/.style={block, fill=blue!12},
  crypto/.style={block, fill=orange!20},
  diffusion/.style={block, fill=red!15},
  arrow/.style={->, >=stealth, thick},
  label/.style={font=\footnotesize\itshape, gray},
]

%% Input
\node[io] (key)   {Khóa $K$ + \texttt{nalu\_index} $n$};
\node[io, below=0.7cm of key] (dc) {DC bytes $d[0..24]$};

%% Step 1: PLCM
\node[proc, below=0.7cm of dc] (plcm)
  {PLCM($x_0$, $p$, $n$) → $ks[0..24]$};

%% Step 2: Arnold
\node[crypto, below=0.9cm of plcm] (arnold)
  {Arnold 2D ($P=3$, $Q=11$, $N=5$)\\ $d[0..24]$ → $d'[0..24]$};

%% Step 3: XOR + Diffusion
\node[diffusion, below=0.9cm of arnold] (xor)
  {Chained XOR:\\ $C[i] = (d'[i] \oplus ks[i]) \oplus C[i-1]$};

%% Output
\node[io, below=0.8cm of xor] (out) {Ciphertext $C[0..24]$};

%% Self-loop for diffusion feedback
\draw[arrow, dashed, red!60] (xor.east) -- ++(1.2,0)
  |- node[right, font=\tiny] {$C[i-1]$} (xor.east |- xor.north);

%% Main arrows
\draw[arrow] (key) -- (plcm) node[midway, right, label] {dẫn xuất $x_0$, $p$};
\draw[arrow] (dc)  -- (plcm);
\draw[arrow] (plcm) -- (arnold) node[midway, right, label] {keystream $ks$};
\draw[arrow] (dc |- arnold.north) -- (arnold)
  node[pos=0, right, label] {};
\draw[arrow] (dc) -- ++(3.5, 0) |- (arnold.east);
\draw[arrow] (arnold) -- (xor);
\draw[arrow] (xor) -- (out);

%% Decrypt note
\node[right=3.8cm of xor, font=\small, align=left, text width=4cm]
  {\textit{Giải mã (ngược):}\\
   $d'[0] = C[0] \oplus ks[0]$\\
   $d'[i] = C[i] \oplus C[i-1] \oplus ks[i]$\\
   Sau đó Arnold$^{-1}$ → $d$};

\end{tikzpicture}

\caption{Sơ đồ luồng thuật toán hybrid: PLCM sinh keystream, Arnold hoán vị,
         XOR + chained diffusion tạo ciphertext; bên phải là sơ đồ giải mã ngược}
\label{fig:hybrid_flowchart}
\end{figure}

\subsection{Bước 1: Sinh keystream PLCM (per-NALU)}

\subsubsection{Hàm dẫn xuất khóa (Key Derivation Function)}

Từ khóa $K$ (chuỗi ký tự ASCII), hàm dẫn xuất khóa tạo ra hai tham số PLCM
$(x_0, p)$:

\begin{enumerate}
  \item Tính hash của $K$: $h = \text{SHA-256}(K)$ (trong thực tế, dùng custom hash đơn giản hơn).
  \item Lấy 8 bytes đầu của $h$, chuyển thành double $x_0' \in [0, 1)$.
  \item Ràng buộc vào vùng hỗn độn: $x_0 = 0.0001 + x_0' \times 0.4998$.
  \item Tương tự cho $p$: $p = 0.0001 + h' \times 0.2499$ (với $h'$ từ bytes 8--15).
\end{enumerate}

Kết quả: $x_0 \in (0.0001, 0.4999)$ và $p \in (0.0001, 0.2500)$ --- đảm bảo
PLCM hoạt động trong vùng hỗn độn cho mọi khóa K.

\subsubsection{Sinh điều kiện đầu per-NALU}

Thay vì dùng trực tiếp $x_0$ cho mọi NALU (keystream cố định), thuật toán
lặp PLCM $n = \texttt{nalu\_index}$ lần để tạo điều kiện đầu riêng cho NALU thứ $n$:

\begin{equation}
x_n = f^n(x_0, p)
\label{eq:plcm_advance}
\end{equation}

Với $f$ là PLCM (Phương trình 2.1). Vì PLCM nhạy cảm với điều kiện đầu,
$x_n$ và $x_{n+1}$ khác nhau hoàn toàn ngay cả khi chỉ cách nhau 1 lần lặp.

\textbf{Tối ưu hoá thực tế:} Thay vì lặp $n$ lần cho mỗi NALU (chi phí $O(n)$),
thực tế cài đặt duy trì trạng thái PLCM liên tục qua các NALU:
\begin{itemize}
  \item NALU đầu tiên: $x_{\text{curr}} = f(x_0, p)$ (1 lần lặp).
  \item NALU tiếp theo: $x_{\text{curr}} = f(x_{\text{curr}}, p)$ (thêm 1 lần lặp).
  \item Chi phí trung bình: O(1) mỗi NALU.
\end{itemize}

\subsubsection{Sinh keystream 25 bytes}

Từ điều kiện đầu $x_n$ của NALU thứ $n$, sinh keystream 25 bytes:

\begin{equation}
ks[i] = \left\lfloor f^{(i)}(x_n, p) \times 256 \right\rfloor \bmod 256, \quad i = 0, 1, \ldots, 24
\end{equation}

Sau đó lấy phần nguyên và modulo 256 để chuyển từ $[0, 1)$ sang byte $[0, 255]$.

\textbf{Ví dụ minh hoạ:} Giả sử $x_n = 0.314159$, $p = 0.1234$:
\begin{itemize}
  \item $f(0.314159, 0.1234) = (0.314159 - 0.1234) / (0.5 - 0.1234) = 0.5064$
        (rơi vào nhánh thứ hai: $x \in [p, 0.5)$)
  \item $ks[0] = \lfloor 0.5064 \times 256 \rfloor = \lfloor 129.64 \rfloor = 129$
\end{itemize}

\subsection{Bước 2: Hoán vị Arnold 2D Cat Map}

\subsubsection{Thuật toán hoán vị}

Mảng DC 25 bytes được xử lý như ma trận $5 \times 5$. Cho mỗi tọa độ $(x, y)$
với $x, y \in \{0, 1, 2, 3, 4\}$:

\begin{equation}
\begin{pmatrix} x' \\ y' \end{pmatrix}
= \begin{pmatrix} 1 & 3 \\ 11 & 34 \end{pmatrix}
\begin{pmatrix} x \\ y \end{pmatrix} \pmod{5}
\end{equation}

Bảng hoán vị đầy đủ với $P=3, Q=11, N=5$:

\begin{table}[ht]
\centering
\caption{Hoán vị Arnold map với $P=3, Q=11, N=5$: chỉ số byte gốc $\to$ chỉ số đích}
\label{tab:arnold_perm}
\begin{tabular}{ccc|ccc}
\toprule
$(x,y)$ gốc & Chỉ số gốc $i$ & Chỉ số đích $i'$ &
$(x,y)$ gốc & Chỉ số gốc $i$ & Chỉ số đích $i'$ \\
\midrule
(0,0) & 0  & 0  & (0,3) & 15 & 15 \\
(1,0) & 1  & 11 & (1,3) & 16 & 1  \\
(2,0) & 2  & 22 & (2,3) & 17 & 12 \\
(3,0) & 3  & 8  & (3,3) & 18 & 23 \\
(4,0) & 4  & 19 & (4,3) & 19 & 9  \\
(0,1) & 5  & 20 & (0,4) & 20 & 5  \\
(1,1) & 6  & 6  & (1,4) & 21 & 16 \\
(2,1) & 7  & 17 & (2,4) & 22 & 2  \\
(3,1) & 8  & 3  & (3,4) & 23 & 13 \\
(4,1) & 9  & 14 & (4,4) & 24 & 24 \\
(0,2) & 10 & 10 & & & \\
(1,2) & 11 & 21 & & & \\
(2,2) & 12 & 7  & & & \\
(3,2) & 13 & 18 & & & \\
(4,2) & 14 & 4  & & & \\
\bottomrule
\end{tabular}
\end{table}

Quan sát: byte 0 và byte 24 là \textbf{fixed points} của hoán vị này
($\pi(0) = 0$ và $\pi(24) = 24$). Đây là đặc tính của Arnold map với tham số
$(3, 11)$ và $N = 5$ --- fixed point là nguyên nhân UACI chưa đạt $33.46\%$
(phân tích ở Chương 6, mục 6.4.5).

\subsubsection{Cài đặt C++}

Hoán vị được tính trước (precomputed) thành bảng tra cứu $\pi[25]$ để tránh
tính lại mỗi lần mã hoá:

\begin{lstlisting}[language=C++, caption={Precomputed Arnold permutation --- \texttt{hybrid\_encryption.cpp}}]
// Precomputed perm[src_index] = dst_index for P=3, Q=11, N=5
static const int ARNOLD_PERM[25] = {
    0, 11, 22,  8, 19,
   20,  6, 17,  3, 14,
   10, 21,  7, 18,  4,
   15,  1, 12, 23,  9,
    5, 16,  2, 13, 24
};

// Apply Arnold permutation: permuted[ARNOLD_PERM[i]] = dc[i]
std::vector<uint8_t> applyArnold(const std::vector<uint8_t>& dc) {
    std::vector<uint8_t> permuted(25);
    for (int i = 0; i < 25; ++i)
        permuted[ARNOLD_PERM[i]] = dc[i];
    return permuted;
}
\end{lstlisting}

\subsection{Bước 3: XOR và Khuếch tán chuỗi (Chained Diffusion)}

\subsubsection{Công thức mã hoá}

Sau khi có $\text{permuted}[0..24]$ và $ks[0..24]$, quá trình mã hoá:

\begin{equation}
C[0] = \text{permuted}[0] \oplus ks[0]
\label{eq:c0}
\end{equation}

\begin{equation}
C[i] = \bigl(\text{permuted}[i] \oplus ks[i]\bigr) \oplus C[i-1], \quad i = 1, 2, \ldots, 24
\label{eq:chained_diffusion}
\end{equation}

Khai triển ra với $i = 2$:
\begin{align}
C[2] &= (\text{permuted}[2] \oplus ks[2]) \oplus C[1] \\
     &= (\text{permuted}[2] \oplus ks[2]) \oplus (\text{permuted}[1] \oplus ks[1]) \oplus C[0] \\
     &= (\text{permuted}[2] \oplus ks[2]) \oplus (\text{permuted}[1] \oplus ks[1]) \oplus (\text{permuted}[0] \oplus ks[0])
\end{align}

Tổng quát:
\begin{equation}
C[i] = \bigoplus_{j=0}^{i} \bigl(\text{permuted}[j] \oplus ks[j]\bigr)
\label{eq:chained_expand}
\end{equation}

Phương trình~(\ref{eq:chained_expand}) cho thấy $C[i]$ phụ thuộc vào \emph{tất cả}
$\text{permuted}[0..i]$ và $ks[0..i]$.

\subsubsection{Phân tích hiệu ứng khuếch tán (Avalanche Analysis)}

\textbf{Kịch bản:} Byte $\text{permuted}[k]$ thay đổi 1 bit (do flip 1 bit DC[j] mà $\pi(j) = k$).
Hiệu ứng lan truyền:

\begin{itemize}
  \item $C[k]$ thay đổi (vì phụ thuộc $\text{permuted}[k]$).
  \item $C[k+1]$ thay đổi (vì phụ thuộc $C[k]$).
  \item $C[k+2]$ thay đổi (vì phụ thuộc $C[k+1]$).
  \item $\vdots$
  \item $C[24]$ thay đổi (vì phụ thuộc $C[23]$).
\end{itemize}

Do đó, khi $k = 0$ (flip byte đầu tiên của DC, tức byte 0 sau hoán vị):
tất cả 25 bytes $C[0..24]$ đều thay đổi $\Rightarrow$ NPCR lý thuyết = 100\%.

Khi $k = 24$ (flip byte cuối): chỉ $C[24]$ thay đổi $\Rightarrow$ NPCR = 4\%.

\textbf{NPCR trung bình lý thuyết:} Trung bình qua 25 vị trí byte:
\begin{equation}
\overline{\text{NPCR}}_{\text{theory}} = \frac{1}{25}\sum_{k=0}^{24} \frac{25 - k}{25} \times 100\%
= \frac{1}{25} \times \frac{(25+1) \times 25/2}{25} \times 100\% = \frac{26}{50} \times 100\% = 52\%
\end{equation}

Tuy nhiên, NPCR được đo bằng flip byte 0 của DC gốc, qua hoán vị Arnold:
byte 0 của DC ($i=0$) sau hoán vị $\pi(0) = 0$, nên đổ vào $\text{permuted}[0] = k = 0$.
Do đó NPCR thực tế $\approx 100\%$ (lý thuyết), kết quả đo được 99.53--99.69\%.

\textbf{Lý giải sai lệch 0.31--0.47\%:}
Do \textbf{byte 24 là fixed point} ($\pi(24) = 24$): $\text{permuted}[24] = \text{DC}[24]$.
Khi flip byte 0, $C[0] \to C[23]$ đều thay đổi nhưng với xác suất thay đổi khác nhau
do XOR với $ks$ có thể tình cờ bù trừ. Sai lệch thực nghiệm từ phân phối DC video cụ thể.

\subsubsection{Cài đặt C++ --- Mã hoá}

\begin{lstlisting}[language=C++, caption={Mã hoá hybrid với chained diffusion --- \texttt{hybrid\_encryption.cpp}}]
std::vector<uint8_t> HybridEncryption::encrypt(
        const std::vector<uint8_t>& dc_data,
        const std::vector<uint8_t>& key,
        int nalu_index) {

    // Step 1: Generate PLCM keystream (per-NALU via nalu_index)
    auto ks = generateKeystream(key, nalu_index, 25);  // 25 bytes

    // Step 2: Arnold permutation
    auto permuted = applyArnold(dc_data);              // 25 bytes permuted

    // Step 3: XOR + chained diffusion
    std::vector<uint8_t> cipher(25);
    cipher[0] = permuted[0] ^ ks[0];
    for (int i = 1; i < 25; ++i)
        cipher[i] = (permuted[i] ^ ks[i]) ^ cipher[i-1];

    return cipher;
}
\end{lstlisting}

\subsection{Giải mã (Decryption)}

\subsubsection{Giải mã chained diffusion}

Giải mã chained diffusion là phép ngược của Phương trình~(\ref{eq:chained_diffusion}):

\begin{equation}
\text{permuted}[0] = C[0] \oplus ks[0]
\label{eq:dec0}
\end{equation}

\begin{equation}
\text{permuted}[i] = C[i] \oplus C[i-1] \oplus ks[i], \quad i = 1, 2, \ldots, 24
\label{eq:dec_chain}
\end{equation}

\textbf{Chứng minh tính đúng đắn:}
\begin{align}
C[i] \oplus C[i-1] \oplus ks[i] &= (\text{permuted}[i] \oplus ks[i] \oplus C[i-1]) \oplus C[i-1] \oplus ks[i] \\
&= \text{permuted}[i] \oplus ks[i] \oplus ks[i] \oplus C[i-1] \oplus C[i-1] \\
&= \text{permuted}[i] \qquad \checkmark
\end{align}

Vì $a \oplus a = 0$ với mọi $a$.

\subsubsection{Giải mã Arnold ngược}

Sau khi có $\text{permuted}[0..24]$, áp dụng hoán vị ngược $\pi^{-1}$:

\begin{equation}
\text{DC}[i] = \text{permuted}[\pi(i)], \quad \forall i \in \{0, \ldots, 24\}
\end{equation}

Hoặc tương đương: $\text{DC}[\pi^{-1}(j)] = \text{permuted}[j]$ với hoán vị ngược
được precomputed từ ma trận $\mathbf{A}^{-1}$.

\subsubsection{Tính xác định và byte-exact}

Vì PLCM là hàm xác định (deterministic) và \texttt{nalu\_index} là tham số công khai
(biết được từ vị trí NALU trong file), quá trình giải mã sinh lại đúng keystream
$ks$ đã dùng khi mã hoá. Kết hợp với giải mã XOR-chain và Arnold inverse hoàn toàn
xác định, toàn bộ quá trình giải mã là \textbf{deterministic và lossless}.

\textbf{Kết quả thực nghiệm:} 8/8 tổ hợp đạt BER = 0 và SHA-256 PASS,
xác nhận byte-exact cho toàn bộ file 219 MB.

\subsubsection{Cài đặt C++ --- Giải mã}

\begin{lstlisting}[language=C++, caption={Giải mã hybrid --- \texttt{hybrid\_encryption.cpp}}]
std::vector<uint8_t> HybridEncryption::decrypt(
        const std::vector<uint8_t>& encrypted_data,
        const std::vector<uint8_t>& key,
        int nalu_index) {

    // Step 1: Regenerate same PLCM keystream (same key, same nalu_index)
    auto ks = generateKeystream(key, nalu_index, 25);

    // Step 2: Inverse chained diffusion XOR
    std::vector<uint8_t> permuted(25);
    permuted[0] = encrypted_data[0] ^ ks[0];
    for (int i = 1; i < 25; ++i)
        permuted[i] = encrypted_data[i] ^ encrypted_data[i-1] ^ ks[i];

    // Step 3: Inverse Arnold permutation
    return inverseArnold(permuted);
}
\end{lstlisting}

\subsection{Interface IEncryption --- Kiến trúc mở rộng}

\subsubsection{Thiết kế abstract interface}

Tất cả 4 thuật toán (hybrid, AES, DES, RC4) đều triển khai qua interface chung:

\begin{lstlisting}[language=C++, caption={IEncryption interface --- \texttt{algorithms/encryption\_interface.h}}]
class IEncryption {
public:
    virtual ~IEncryption() = default;

    // Encrypt dc_data (25 bytes) với key và nalu_index
    // nalu_index được dùng bởi hybrid để tạo per-NALU keystream
    // AES/DES/RC4 bỏ qua nalu_index (keystream cố định)
    virtual std::vector<uint8_t> encrypt(
        const std::vector<uint8_t>& dc_data,
        const std::vector<uint8_t>& key,
        int nalu_index) = 0;

    virtual std::vector<uint8_t> decrypt(
        const std::vector<uint8_t>& encrypted_data,
        const std::vector<uint8_t>& key,
        int nalu_index) = 0;

    // Factory method: tạo instance theo tên thuật toán
    // Hỗ trợ: "hybrid", "aes", "des", "rc4"
    static std::unique_ptr<IEncryption> create(const std::string& algo_name);
};
\end{lstlisting}

\subsubsection{Factory Pattern}

\begin{lstlisting}[language=C++, caption={Factory method --- \texttt{algorithms/encryption\_factory.cpp}}]
std::unique_ptr<IEncryption> IEncryption::create(const std::string& algo) {
    if (algo == "hybrid") return std::make_unique<HybridAlgo>();
    if (algo == "aes")    return std::make_unique<AesAlgo>();
    if (algo == "des")    return std::make_unique<DesAlgo>();
    if (algo == "rc4")    return std::make_unique<Rc4Algo>();
    throw std::invalid_argument("Unknown algorithm: " + algo);
}
\end{lstlisting}

Pipeline chỉ gọi \texttt{IEncryption::create(algo\_name)} một lần, sau đó
gọi \texttt{encrypt()} / \texttt{decrypt()} trong vòng lặp NALU. Thêm thuật toán
mới chỉ cần: (1) tạo class mới kế thừa IEncryption, (2) đăng ký trong factory.
Pipeline không cần sửa.

\subsection{Phân tích keyspace}

\subsubsection{Keyspace từng thành phần}

\begin{table}[ht]
\centering
\caption{Phân tích keyspace thuật toán hybrid}
\label{tab:keyspace}
\begin{tabular}{llll}
\toprule
Thành phần & Tham số & Keyspace & Ghi chú \\
\midrule
PLCM điều kiện đầu & $x_0 \in (0.0001, 0.4999)$ & $\approx 2^{52}$ & Giới hạn mantissa 52-bit \\
PLCM tham số & $p \in (0.0001, 0.2500)$ & $\approx 2^{51}$ & Cùng giới hạn \\
PLCM tổng & $(x_0, p)$ & $\approx 2^{53}$ & Kết hợp (overlap do cùng nguồn) \\
Arnold Cat Map & $(P, Q)$ với $N=5$ & $\approx 2^{4}$ & $P, Q \in \{1,2,3,4\}$: 16 cặp \\
Arnold iterations & số lần lặp & $\approx 2^{1}$ & Thực tế cố định = 1 \\
\midrule
\textbf{Tổng ước tính} & & $\approx 2^{58}$ & $2^{53} + 2^{4} \approx 2^{53}$ \\
\bottomrule
\end{tabular}
\end{table}

\subsubsection{Đánh giá và hạn chế}

Keyspace $2^{58}$ tương đương khoảng $2.88 \times 10^{17}$ khóa phân biệt.
Với phần cứng hiện đại (GPU có thể thử $10^{12}$ khóa/giây):
thời gian brute-force $\approx 2^{58} / 10^{12} \approx 288.000$ giây $\approx 3.3$ ngày.

Đây là \textbf{không đủ} cho bảo mật hiện đại (ngưỡng $2^{128}$). Hạn chế này
xuất phát từ việc sử dụng IEEE 754 double-precision (52-bit mantissa) cho PLCM.

Giải pháp đề xuất trong Chương 7: thay PLCM double bằng PLCM multi-precision
128-bit fixed-point, tăng keyspace lên $2^{128}$.

\subsection{Tóm tắt thuật toán}

\begin{table}[ht]
\centering
\caption{Tóm tắt thuật toán hybrid mã hoá và giải mã}
\label{tab:algo_summary}
\begin{tabular}{lll}
\toprule
Bước & Mã hoá & Giải mã \\
\midrule
1 & PLCM: $(K, \texttt{nalu\_index}) \to ks[0..24]$ & PLCM: $(K, \texttt{nalu\_index}) \to ks[0..24]$ \\
2 & Arnold: $\text{DC}[0..24] \to \text{permuted}[0..24]$ & XOR-chain ngược: $C[0..24] \to \text{permuted}[0..24]$ \\
3 & XOR-chain: $\text{permuted} + ks \to C[0..24]$ & Arnold$^{-1}$: $\text{permuted}[0..24] \to \text{DC}[0..24]$ \\
\bottomrule
\end{tabular}
\end{table}
